The Erd\H os--Kac theorem is a corollary of the Kubilius model, and it is already a theorem of Volume~6. One proof is enough to cite. This section says what the corollary is, in two pages, and what each proof makes visible. It does not found the theorem a second time. The Sathe--Selberg formula, which sharpens the count of integers with $\omega(n)=k$, is Cauchy's integral of the Selberg--Delange mean already proved in this chapter; it is not rewritten here.

\begin{theorem}[Erd\H os--Kac, as a corollary]\label{thm:erdoskac-cor}
For every real $a$, the proportion of $n\le x$ with $(\omega(n)-\log\log x)/\sqrt{\log\log x}\le a$ tends to $\Phi(a)$.
\end{theorem}
\begin{proof}
Let $L=\log\log x$ and $y=\exp(L^{1/2})$, so $\log\log y=\tfrac12 L\to\infty$ and $u=\log x/\log y=L^{1/2}\to\infty$. Theorem~\ref{thm:kubilius} gives total variation $o(1)$ between the law of $(\nu_p(n))_{p\le y}$ and the independent geometrics $G_p$. The additive function $\omega_y(n)=\sum_{p\le y}\mathbf 1_{p\mid n}$ is a function of that vector. Under the model, the sum of independent Bernoulli random variables of means $1/p$ has mean and variance $\sum_{p\le y}1/p=\log\log y+O(1)=\tfrac12 L+O(1)$, and the classical central limit theorem, which Volume~4 proves for uniformly bounded independent summands with variance tending to infinity, yields a Gaussian after normalisation by $\sqrt{\log\log y}$.

The primes between $y$ and $x$ change $\omega$ by $\omega(n)-\omega_y(n)\le\log x/\log y=L^{1/2}$. Tur\'an--Kubilius, already proved in this chapter, bounds the second moment of $\omega-\omega_y$ on $n\le x$ by $O(L^{1/2})$. Chebyshev's inequality makes $(\omega-\omega_y)/\sqrt L\to0$ in probability. The means $\log\log x$ and $\log\log y$ differ by $\tfrac12 L$, which is the same order as the variance; replace $y$ by $x^{1/\log L}$, so that $\log\log y=L-\log\log L$ and $u=\log L\to\infty$, which is still inside the range of Theorem~\ref{thm:kubilius}. Then the truncated mean is $L+O(\log\log L)$, the Gaussian normalisation differs from $\sqrt L$ by $o(\sqrt L)$, and adding a quantity that is $o_P(\sqrt L)$ does not change the limiting distribution function. This is the corollary.
\end{proof}

Volume~6 proves the same limit by the method of moments, truncating at $z=x^{1/L^{1/4}}$ and comparing the moments of $\omega_z$ with those of the independent sum, then smoothing. The two arguments use the same truncation. They do not use the same bookkeeping.

What Volume~6 makes visible is the moment comparison: after the large prime factors are cut away, $\omega$ has the moments of a sum of independent coin tosses, and those moments are the Gaussian moments. The proof never builds a probability space on the primes; the independence is an asymptotic statement about averages of $\mathbf 1_{p\mid n}$.

What this chapter makes visible is the space. Theorem~\ref{thm:kubilius} is a total-variation comparison, so the coin tosses are a measure, not only a list of moments. The bell curve is the central limit theorem on that product, which Volume~4 already contains. The arithmetic content is the estimate that says the product is close to the integers.

The rate in the corollary is not the one R\'enyi and Tur\'an obtained from the Selberg--Delange formula. That rate is $O(1/\sqrt{\log\log x})$ in the Kolmogorov distance, and it uses the mean of $z^{\omega(n)}$ for complex $z$ of modulus $1$, which this chapter has proved and which Volume~6 does not need. Recording the rate would be a second pass through the same generating function, not a second foundation. Volume~11's description of a uniform integer modulo $\varphi$, which quotes the Selberg--Delange theorem from this chapter, is a consumer of the mean-value theorem, not of a second proof of Erd\H os--Kac.

\begin{proposition}[a Chernoff bound]\label{prop:chernoffomega}
Let $Q^*(\alpha)=\alpha\log\alpha-\alpha+1$ for $\alpha>0$, and fix $A\ge1$. Uniformly for $x\ge16$,
\[
P_x\bigl(\omega(n)\ge\alpha\log\log x\bigr)\ll_A(\log x)^{-Q^*(\alpha)}\qquad(1\le\alpha\le A).
\]
\end{proposition}
\begin{proof}
Selberg's formula in this chapter gives $\mathbb E_x z^{\omega(n)}\ll_A(\log x)^{z-1}$ for real $z\in[1,A]$. Markov's inequality at $z=\alpha$ produces the exponent $\alpha-1-\alpha\log\alpha=-Q^*(\alpha)$.
\end{proof}
